Una massa de 100~g es fa oscil·lar penjada d'una molla. S'observa que fa 40 oscil·lacions en un minut i que la diferència entre la posició més alta i la més baixa és de 15~cm.

\begin{apartat}{1,25}
Determineu el període, la constant de la molla i l'equació del moviment si comencem a comptar el moviment quan passa per la posició més baixa. Representeu a la quadrícula de sota la força elàstica durant dos períodes sencers.
\end{apartat}

\figura[0.6]{fig1.png}

\begin{apartat}{1,25}
Calculeu l'energia mecànica de l'oscil·lador harmònic i trobeu l'expressió de l'energia cinètica en funció de la posició de la massa. Calculeu el mòdul de la velocitat quan la massa és 3~cm per sobre de la posició d'equilibri.
\end{apartat}
